\documentclass[10pt,a4paper]{amsart}
\usepackage[margin=19mm]{geometry}
\usepackage{amsmath,amssymb,mathtools}
\usepackage[scr=rsfs,cal=euler]{mathalfa}
\usepackage{xfrac}
\usepackage[unicode,hidelinks,bookmarks=false]{hyperref}
\newcommand{\ie}{i.e.\ }
\newcommand{\cf}{cf.\ }
\newcommand{\resp}{respectively\ }
\newcommand{\Subsection}{\subsection}
\providecommand{\nobreakdash}{\nobreak\hbox{-}}
\setlength{\emergencystretch}{2em}
\multlinegap0pt
\usepackage{bm}
\usepackage{bbm}
\usepackage{dsfont}
\usepackage{xspace}
\usepackage{cancel}
\usepackage{mathtools}
\usepackage{amsthm}
\makeatletter
\@ifundefined{thm}{\newtheorem{thm}{Theorem}}{}
\@ifundefined{cor}{\newtheorem{cor}[thm]{Corollary}}{}
\@ifundefined{lem}{\newtheorem{lem}[thm]{Lemma}}{}
\@ifundefined{prop}{\newtheorem{prop}[thm]{Proposition}}{}
\makeatother
\newcommand{\I}{\mathbb{I}}
\newcommand{\T}{\mathbf{T}}
\newcommand{\Varpi}{\rotatebox[origin=c]{180}{$\Pi\kern-0.361em\Pi$}}
\newcommand{\f}{\mathbf{f}}
\newcommand{\ip}{\I_P}
\newcommand{\n}{\mathfrak{n}}
\newcommand{\p}{\mathfrak{p}}
\newcommand{\qp}{\mathbb{Q}_{p}}
\newcommand{\tp}{\widetilde{P}}
\newcommand{\z}{\mathbb{Z}}
\title[A $p$-Converse theorem for Real Quadratic Fields]{A $p$-Converse theorem for Real Quadratic Fields}
\author{Muskan Bansal, Somnath Jha, Aprameyo Pal, Guhan Venkat}
\date{Source: arXiv:2504.21799v3 (CC BY 4.0)}
\begin{document}
\maketitle

\noindent\emph{Source and licence.} A $p$-Converse theorem for Real Quadratic Fields. Version arXiv:2504.21799v3 (CC BY 4.0), distributed under the Creative Commons Attribution 4.0 International licence, as recorded in the arXiv licence metadata for this exact version and shown on the abstract page https://arxiv.org/abs/2504.21799v3. Full rights notice: licenses/arxiv-2504.21799-CC-BY-4.0.txt.\par\smallskip

\noindent\textit{Editorial scope.} Counted unit: the Thm whose source label is \texttt{None} in the article (source0.tex lines 1503--1514, source label \texttt{None}), delivered together with the article's own complete original proof of it (source0.tex lines 1516--1560, one \texttt{proof} environment of 45 lines). No mathematics is written, repaired or re-derived: statement and proof are the article's own text line for line. Required context retained uncounted: ses for MW group (lines 846--848); iedagger is an iso (lines 1039--1041); NonDegeneracyOfPairing (lines 1096--1102); NonDegeneracyEquivalence (lines 1278--1289); PairingInTermsOfLog (lines 1435--1445). Every definition, display and label of those lines is retained unchanged. Omitted from this package and not redistributed: 1--845, 849--1038, 1042--1095, 1103--1277, 1290--1434, 1446--1502, 1515--1515, 1561--1895. Nothing inside a retained excerpt is omitted or altered: every retained body is delivered exactly as the article's lines are written, so no omission receipt of any kind accompanies this package. Citations: the retained excerpts carry no citation command at all, so no bibliography entry of the article is transcribed and no reference is redistributed. Reference adaptation: none was needed and none was made; every reference inside the delivered unit resolves inside it, so no unresolved reference or citation is produced. Printed slips kept literal: no printed word, symbol, hypothesis or number of the counted statement or of its proof was altered. Layout and class: the article's own class is \texttt{amsart} with packages including fontenc, graphicx, mathtools, amssymb, amsthm, thmtools, mathrsfs, euscript, bigints, xcolor, tikz-cd, nameref, dsfont, hyperref; the wrapper is the lane's pinned \texttt{amsart} editorial wrapper at 10pt on A4, which restates the theorem-like environments the retained text needs and adds the article's own macro definitions that the retained text needs (\texttt{\textbackslash I}, \texttt{\textbackslash T}, \texttt{\textbackslash Varpi}, \texttt{\textbackslash f}, \texttt{\textbackslash ip}, \texttt{\textbackslash n}, \texttt{\textbackslash p}, \texttt{\textbackslash qp}, \texttt{\textbackslash tp}, \texttt{\textbackslash z}), with extra wrapper packages (bm, bbm, dsfont, xspace, cancel, mathtools). The delivered numbering is the unit's own. Assets: no retained span contains a figure environment, a graphics-inclusion command, a PDF-page inclusion, a table or a TikZ picture; no image file is copied into this package, the package declares no asset dependency and no image file is redistributed with it. Licence: version arXiv:2504.21799v3 of this article is distributed under the Creative Commons Attribution 4.0 International licence, as recorded in the arXiv licence metadata for this exact version and shown on the abstract page https://arxiv.org/abs/2504.21799v3. A full rights notice is delivered at licenses/arxiv-2504.21799-CC-BY-4.0.txt. Characters: every retained excerpt is pure ASCII, so the delivered engine, XeLaTeX with its default Latin Modern text font, reports no missing character.
\begin{align} \label{ses for MW group}
    0 \rightarrow E(K') \otimes L \rightarrow Sel_E(K') \otimes L \rightarrow T_p(\Varpi(E/K')) \otimes L \rightarrow 0,
\end{align}

\begin{lem} \label{iedagger is an iso}
    Consider the map $i_E^\dagger : E^\dagger(K') \otimes L \rightarrow \widetilde{H}^1_f(K',V_f)$, induced by $i_E^\dagger$. Then $i_E^\dagger$ is injective. Moreover, if we assume that $\Varpi(E/K')_{p^\infty}$ is finite, then $i_E^\dagger$ is an isomorphism.
\end{lem}

\begin{prop} \label{NonDegeneracyOfPairing}
    \begin{enumerate}
        \item $\langle x,x \rangle^{Nek}_{K'} = 0$ for every $x \in \widetilde{H}^1_f(K',V_f)$.

        \item $\langle -,- \rangle^{Nek}_{K'}$ is a $Gal(K'/F)$-equivariant pairing.
    \end{enumerate}
\end{prop}

\begin{cor} \label{NonDegeneracyEquivalence}
    Let $\chi \in \{1, \epsilon_K\}$ be a quadratic character of $F$ with conductor coprime to $\n\p$ such that $\chi(\p) = 1$. Then  the following are equivalent
    \begin{enumerate}
        \item $\langle -,- \rangle^{Nek}_{K',\chi}$ is non-degenerate.
        \item $\widetilde{H}^2_f(K',\T_{\f,P})^\chi$ is a torsion, semi-simple $\ip$-module.
        \item $\mathrm{len}_P(\widetilde{H}^2_f(K',\T_{\f,P})^\chi) = \mathrm{dim}_L(\widetilde{H}^1_f(K',V_f)^\chi)$.
    \end{enumerate}
    Moreover, if this is the case then, $X_{Gr}^{cc}(\f/K')^\chi_P$ is also a torsion, semi-simple $\ip$-module and
    \begin{align*}
        \mathrm{len}_P(X^{cc}_{Gr}(\f/K')^\chi) = \mathrm{dim}_L(H^1_f(K',V_f)^\chi).
    \end{align*}
\end{cor}

\begin{cor} \label{PairingInTermsOfLog}
    For every $(P,\widetilde{P}) \in E^\dagger(K')$ and $w|\p$, we have
    \begin{align*}
        \langle q_w,(P,\widetilde{P}) \rangle^{Nek}_{K'} = -\frac{1}{2} \log_{N(q_E)}(N_{K'_w/\qp}(\widetilde{P}_w)),
    \end{align*}
    where, as earlier, $N := N_{F_\p/\qp}$ is the usual field norm and $q_E \in F_\p^\times$ is the Tate period of $E/F_\p$. In particular, for $i,j = 1,2$
    \begin{align*}
        \langle q_{w_i},q_{w_j} \rangle^{Nek}_{K'} = 0
    \end{align*}
    where $w_1$ and $w_2$ are primes dividing $\p$.
\end{cor}

\begin{thm}
 Let $\chi \in \{1, \epsilon_K\}$ be a quadratic character of $F$ with conductor coprime with $\mathfrak{np}$. Assume that 
 \begin{enumerate}
     \item $\chi(\p) = 1$,
     \item $\mathrm{rank}_\z \hspace{0.01mm} \hspace{1mm} E(K')^\chi = 1$,
     \item $\Varpi(E/K')^\chi_{p^\infty}$ is finite.
 \end{enumerate}
 Then
 \begin{align*}
     X^{cc}_{Gr}(\f/K')^\chi \otimes_\I \I_P \cong \I_P/PI_P.
 \end{align*}
\end{thm}

\begin{proof}
    Let $P \in E(K')^\chi$ be a generator of the free part of $E(K')^\chi$. Fix a lift $P^\dagger = (P,(\tp_w)_{w|\p}) \in E^\dagger(K')^\chi$ of $P$. If $\chi$ is the trivial character, set
    \begin{align*}
        q_\chi := (0,q_E) \in E^\dagger(F) \subset E(F) \times F_\p^\times,
    \end{align*}
    otherwise, set
    \begin{align*}
        q_\chi := (0,(q_E,q_E^{-1})) \in E^\dagger(K)^\chi \subset E(K) \times K_{w_1}^\times \times K_{w_2}^\times,
    \end{align*}
    where $\p O_K = w_1w_2$. Since $\Varpi(E/K')^{\chi}_{p^\infty}$ is finite, Lemma \ref{iedagger is an iso} and the hypothesis $\mathrm{rank}_\z \hspace{0.01mm} \hspace{1mm} E(K')^\chi = 1$ imply that 
    \begin{align*}
        \widetilde{H}^1_f(K',V_f)^\chi \cong (E^\dagger(K') \otimes L)^\chi \cong L \cdot P^\dagger \oplus L \cdot q_\chi.
    \end{align*}
    By $\doteq$, we mean equality up to some nonzero constant multiple. If $K'=F$, using Corollary \ref{PairingInTermsOfLog} we have the following equality 
    \begin{align*}
        \langle q_\chi,P^\dagger \rangle^{Nek}_F \doteq \log_{E,N}(P),
    \end{align*}
    where $\log_{E,N}(\cdot) := \log_{N(q_E)} \circ N \circ \Phi_{Tate}^{-1}(\cdot)$.
    \\
    If $K' = K$ is a quadratic extension of $F$, since $\p$ splits in $K$, write $i_1 : K \hookrightarrow K_{w_1} \cong F_\p$ and $i_2 : K \hookrightarrow K_{w_2} \cong F_\p$. Then $i_2 = i_1 \circ \sigma$, where $\sigma \in Gal(K/F)$ is the non trivial element. Since $P^\dagger \in E^\dagger(K)^\chi$, we have $P^\sigma = -P$ and $\tp_{w_1} = \tp_{w_2}^{-1}$.
    Now clearly, $q_\chi = q_{w_1}-q_{w_2}$, we get again using Corollary \ref{PairingInTermsOfLog} and the fact that $K_{w_j} \cong F_\p$ for both $j = 1,2$,
    \begin{align*}
        \langle q_\chi, P^\dagger \rangle^{Nek}_K & \doteq \log_{N(q_E)}(N_{K_{w_1}/\qp}(\tp_{w_1})) - \log_{N(q_E)}(N_{K_{w_2}/\qp}(\tp_{w_2})) \\
        & \doteq \log_{E,N}(i_1(P)) - log_{E,N}(i_2(P)) \\
        & \doteq 2 \log_{E,N}(P).
    \end{align*}
    Note that these equalities hold upto multiplication by a nonzero constant. Since $\widetilde{H}^1_f(K',V_f)^\chi$ is generated by $P^\dagger$ and $q_\chi$, using Proposition \ref{NonDegeneracyOfPairing} we obtain
    \begin{align*}
        det \langle -,- \rangle^{Nek}_{K,\chi} & = det \begin{pmatrix}
            \langle q_\chi, q_\chi \rangle^{Nek}_K & \langle q_\chi, P^\dagger \rangle^{Nek}_K \\
            \langle P^\dagger,q_\chi \rangle^{Nek}_K & \langle P^\dagger, P^\dagger \rangle^{Nek}_K
        \end{pmatrix} \\
        & \doteq det \begin{pmatrix}
            0 & \log_{E,N}(P) \\
            -\log_{E,N}(P) & 0
        \end{pmatrix} \\
        & \doteq \log_{E,N}^2(P).
    \end{align*}
    Using the fact that the only zeros of $\log_{N(q_E)}$ are the roots of unity and powers of $N(q_E)$ and the description of $N$ in degree 2 extension, we conclude that since $P$ is of infinite order, $\log_{E,N}(P)$ is nonzero and hence so is $det \langle -,- \rangle^{Nek}_{K',\chi}$. In particular, $\langle -,- \rangle^{Nek}_{K',\chi}$ is non-degenerate.
    Hence, using Corollary \ref{NonDegeneracyEquivalence} and the short exact sequence (\ref{ses for MW group}), we get that
    \begin{align*}
        \mathrm{len}_P(X^{cc}_{Gr}(\f/K')^\chi) = \mathrm{dim}_L H^1_f(K',V_f)^\chi = 1.
    \end{align*}
    This implies that $X^{cc}_{Gr}(\f/K')^\chi \otimes \I_P \cong \I_P/P\I_P$.
\end{proof}

\end{document}
